\chapter{Chapter 9: Complete Full-Length Mock Examinations}

\section{Overview}
This chapter provides two complete, full-length university examination papers modeled on the official LPU examination pattern. Each paper contains Section A (Objective/MCQs), Section B (Short Answer \& Derivation Problems), and Section C (Long Construction \& AutoCAD Problems) with complete marking rubrics, answer keys, and step-by-step solutions.

---

\section{MEC136 MOCK EXAMINATION PAPER 1}

\noindent \textbf{Time Allowed:} 3 Hours $\qquad \bullet \qquad$ \textbf{Maximum Marks:} 100 Marks

\subsection{SECTION A: Objective \& Multiple-Choice Questions (20 Marks)}
\noindent *Answer all 10 questions. Each question carries 2 Marks.*

1. What is the $R.F.$ of a scale if $1\text{ cm}$ represents $5\text{ m}$?  
   (A) $1:5$ $\quad$ (B) $1:50$ $\quad$ (C) $1:500$ $\quad$ (D) $1:5000$

2. In First Angle Projection, the Top View is placed:  
   (A) Above Front View $\quad$ (B) Below Front View $\quad$ (C) To the left of Front View $\quad$ (D) On $XY$ line

3. The isometric length of an edge whose true length is $100\text{ mm}$ is:  
   (A) $81.6\text{ mm}$ $\quad$ (B) $86.6\text{ mm}$ $\quad$ (C) $70.7\text{ mm}$ $\quad$ (D) $100\text{ mm}$

4. Which line type is used for center lines and axes of symmetry?  
   (A) Continuous Thick $\quad$ (B) Dashed Medium $\quad$ (C) Chain Thin $\quad$ (D) Continuous Thin

5. What is the sector angle $\theta$ for developing a cone of base radius $20\text{ mm}$ and slant height $80\text{ mm}$?  
   (A) $45^\circ$ $\quad$ (B) $90^\circ$ $\quad$ (C) $120^\circ$ $\quad$ (D) $180^\circ$

6. In AutoCAD, which command sets the invisible drawing area boundaries?  
   (A) \texttt{UNITS} $\quad$ (B) \texttt{LIMITS} $\quad$ (C) \texttt{ZOOM} $\quad$ (D) \texttt{GRID}

7. Which solid component is NOT hatched when cut longitudinally?  
   (A) Hollow Pipe $\quad$ (B) Thin Reinforcing Web $\quad$ (C) Casing Wall $\quad$ (D) Flange Rim

8. Which AutoCAD modifier command requires a green Crossing Window selection?  
   (A) \texttt{TRIM} $\quad$ (B) \texttt{EXTEND} $\quad$ (C) \texttt{STRETCH} $\quad$ (D) \texttt{MIRROR}

9. The angle between any two principal isometric axes is:  
   (A) $30^\circ$ $\quad$ (B) $60^\circ$ $\quad$ (C) $90^\circ$ $\quad$ (D) $120^\circ$

10. Which development method is used for Pyramids and Cones?  
    (A) Parallel Line Method $\quad$ (B) Radial Line Method $\quad$ (C) Triangulation Method $\quad$ (D) Box Method

---

\subsection{SECTION B: Short Answer \& Analytical Problems (30 Marks)}
\noindent *Answer any 3 questions. Each question carries 10 Marks.*

11. Derive the mathematical reduction factor ($0.816$) for an Isometric Scale from first principles of trigonometry.
12. A line $AB$, $70\text{ mm}$ long, has its end $A$ $15\text{ mm}$ above HP and $20\text{ mm}$ in front of VP. The line is parallel to VP and inclined at $45^\circ$ to HP. Draw its Front View and Top View, and calculate its apparent top view length.
13. Differentiate between Full Section, Half Section, and Offset Section with neat sketches. State two hatching exception rules.
14. Explain the function of AutoCAD drawing aids: \texttt{OSNAP} (\texttt{F3}), \texttt{ORTHO} (\texttt{F8}), \texttt{LTSCALE}, and \texttt{PRESSPULL}.

---

\subsection{SECTION C: Long Construction \& CAD Workflow Problems (50 Marks)}
\noindent *Answer any 2 questions. Each question carries 25 Marks.*

15. **Line Projection (25 Marks):**  
    A line $PQ$, $85\text{ mm}$ long, has its end $P$ $20\text{ mm}$ above HP and $15\text{ mm}$ in front of VP. The front view of the line measures $65\text{ mm}$ and the top view measures $70\text{ mm}$. Draw the orthographic projections of line $PQ$ using the Rotating Line Method. Determine: (a) True inclination with HP ($\theta$), (b) True inclination with VP ($\phi$), and (c) Locate its Horizontal Trace ($\text{HT}$) and Vertical Trace ($\text{VT}$).

16. **Surface Development (25 Marks):**  
    A square prism of base side $40\text{ mm}$ and height $65\text{ mm}$ rests vertically on its base on HP with one base edge inclined at $30^\circ$ to VP. It is cut by a section plane inclined at $45^\circ$ to HP, perpendicular to VP, and bisecting the vertical axis. Draw: (a) Sectional Front View, (b) Top View, and (c) Full development of lateral surfaces using the Parallel Line Method.

17. **AutoCAD 3D \& 2D Extraction (25 Marks):**  
    Detail the complete step-by-step AutoCAD command sequence and workflow to model a 3D composite solid consisting of a square pyramid (base $40\text{ mm}$, height $50\text{ mm}$) resting coaxially on top of a cylinder (diameter $60\text{ mm}$, height $40\text{ mm}$). Explain how to extract 2D orthographic views onto a Layout sheet using \texttt{VIEWBASE}.

---

\section{ANSWER KEY \& DETAILED SOLUTIONS — PAPER 1}

\subsection{Section A Answer Key}
1. **C ($1:500$)** — $R.F. = 1\text{ cm} / 500\text{ cm} = 1/500$.
2. **B (Below Front View)** — First angle projection layout.
3. **A ($81.6\text{ mm}$)** — Isometric length $= 0.816 \times 100 = 81.6\text{ mm}$.
4. **C (Chain Thin)** — Long-short-long chain line for center lines.
5. **B ($90^\circ$)** — $\theta = (r/L) \times 360^\circ = (20/80) \times 360^\circ = 90^\circ$.
6. **B (\texttt{LIMITS})** — Sets paper boundaries.
7. **B (Thin Reinforcing Web)** — Longitudinal web cuts are left unhatched.
8. **C (\texttt{STRETCH})** — Requires green Crossing Window selection.
9. **D ($120^\circ$)** — Isometric axes meet at $120^\circ$.
10. **B (Radial Line Method)** — Used for pyramids and cones.

\subsection{Section B \& C Key Solutions}

\noindent \textbf{Solution to Question 11 (Isometric Scale Factor Derivation):}  
Consider a square top face of a cube tilted at $45^\circ$ to the true horizontal line and $30^\circ$ to the isometric axis.
$$\text{Isometric Length} = B D \cos 45^\circ, \quad \text{True Length} = B D \cos 30^\circ$$
$$\frac{\text{Isometric Length}}{\text{True Length}} = \frac{\cos 45^\circ}{\cos 30^\circ} = \frac{1/\sqrt{2}}{\sqrt{3}/2} = \sqrt{\frac{2}{3}} \approx \mathbf{0.8165}$$

\noindent \textbf{Solution to Question 12 (Line Parallel to VP):}  
Front View length $= 70\text{ mm}$ at $45^\circ$ to $XY$. Top View is horizontal length $= 70 \cos 45^\circ = 70 \times 0.7071 = \mathbf{49.5\text{ mm}}$.

\noindent \textbf{Solution to Question 15 (Line Inclined to Both Planes):}  
1. $TL = 85\text{ mm}$, $a'b' = 65\text{ mm}$, $ab = 70\text{ mm}$.
2. True inclination with HP: $\cos \theta = \frac{\text{Top View Length}}{TL} = \frac{70}{85} = 0.8235 \implies \mathbf{\theta = 34.56^\circ}$.
3. True inclination with VP: $\cos \phi = \frac{\text{Front View Length}}{TL} = \frac{65}{85} = 0.7647 \implies \mathbf{\phi = 40.12^\circ}$.
